% IEEE conference-format manuscript. Compile: pdflatex paper.tex (twice). Requires IEEEtran.cls (standard in TeX Live / Overleaf).
% The bibliography is generated from build_refs.py (single source of truth, also writes references.bib).
\documentclass[conference]{IEEEtran}
\usepackage{graphicx}
\usepackage{booktabs}
\usepackage{amsmath}
\usepackage{cite}
\usepackage{url}
\usepackage[T1]{fontenc}
\graphicspath{{../../results/figures/}}
\newcommand{\tabsize}{\footnotesize}

\begin{document}

\title{Does Dimensionality Reduction Help Classical Intrusion Detection? Original Features versus PCA and Kernel PCA on a Small CIC-IDS2017 Sample}

\author{\IEEEauthorblockN{[Author One]}
\IEEEauthorblockA{[Department], [Institution]\\
{[City, Country]}\\
{[email]}}
\and
\IEEEauthorblockN{[Author Two]}
\IEEEauthorblockA{[Department], [Institution]\\
{[City, Country]}\\
{[email]}}}

\maketitle

\begin{abstract}
Kernel PCA (KPCA) is often paired with classical classifiers for network intrusion detection, but reported gains are usually measured with accuracy on a single split and without a no-reduction baseline. We compare the original 78 flow features, PCA, and KPCA (linear, RBF, cosine; 5--30 components) across five classifiers (logistic regression, $k$-NN, SVM, decision tree, random forest) on a 1{,}000-record stratified sample of CIC-IDS2017 (11 classes, 71.8\% benign). Scaling and reduction are fitted on training folds only; the cached sample carries a residual pre-sampling imputation leak, disclosed below. We use repeated stratified 5-fold cross-validation (15 folds), nested cross-validation for configuration selection, macro-F1, balanced accuracy and MCC, Wilcoxon tests with Holm correction, and a runtime/memory benchmark. The highest observed macro-F1 among the evaluated configurations was a random forest on the \emph{original} features (accuracy $0.952\pm0.010$, macro-F1 $0.722\pm0.058$; nested estimate $0.715\pm0.044$), while web-attack F1 was only 0.21--0.50, so accuracy overstates detection quality. Original features generally outperformed PCA and KPCA for decision tree and random forest (clearly) and logistic regression (modestly); $k$-NN was broadly unchanged; the main exception was an RBF-SVM helped by 30-component cosine KPCA (macro-F1 0.39 to 0.51). Linear KPCA reproduced PCA. KPCA raised fitting time roughly tenfold and traced memory from 1.7 to 11.7\,MB at 800 training records, growing roughly quadratically (a dense kernel matrix for all 1.3\,M flows would theoretically need about 14\,TB; an arithmetic estimate, not a measurement). Results come from one small sample of one dataset and do not establish generalization.
\end{abstract}

\begin{IEEEkeywords}
network intrusion detection, kernel PCA, dimensionality reduction, CIC-IDS2017, class imbalance, empirical comparison
\end{IEEEkeywords}

\section{Introduction}
Machine-learning network intrusion detection systems (NIDS) are commonly developed and compared on public flow datasets \cite{ring2019}, of which CIC-IDS2017 \cite{sharafaldin2018} is widely used. Its flow records carry 78 numeric features, which motivates dimensionality reduction as a way to cut training cost and possibly improve generalization. Linear PCA and auto-encoders have been applied to CIC-IDS2017 \cite{abdulhammed2019}, and Kernel PCA (KPCA) \cite{scholkopf1998}, which performs PCA in an implicit nonlinear feature space, has been combined with SVMs for intrusion detection \cite{kuang2014}.

Three methodological issues make such comparisons hard to interpret. First, CIC-IDS2017 is class-imbalanced and, in its original form, contains labelling and flow-construction problems \cite{engelen2021,panigrahi2018}, so accuracy can be dominated by the benign majority. Second, pitfalls in security-oriented machine learning can lead to over-optimistic performance estimates \cite{arp2022}. Third, a no-reduction baseline is the necessary control for any claim that reduction helps; without it, a benefit is assumed rather than measured.

\textbf{Research questions.} We ask: \textbf{RQ1} how does reduction affect multi-class detection performance; \textbf{RQ2} does KPCA beat the original space or conventional PCA; \textbf{RQ3} how do kernel, component count and classifier interact; \textbf{RQ4} what does the reduction cost in time and memory; \textbf{RQ5} how stable are the conclusions under repeated, class-aware validation and smaller training sets.

\textbf{Contributions.} (i) A reproducible comparison of original features, PCA and three KPCA kernels over five classical classifiers with fold-wise fitting of scaling and reduction, a majority-class baseline and a nested-CV check of selection optimism (the cached sample retains a disclosed residual imputation leak). (ii) A class-aware analysis (macro-F1, balanced accuracy, MCC, per-class recall) showing how much accuracy overstates detection quality on this data. (iii) Paired statistical tests of the comparisons that matter, with an explicit caveat on their independence. (iv) A runtime and memory measurement of KPCA against PCA. We do not propose a new algorithm. All code, raw per-fold results and scripts are in the project repository.

\section{Related Work}
\textit{Datasets.} Ring et al.\ survey network-based intrusion datasets and the properties that determine their suitability for evaluation scenarios \cite{ring2019}. CIC-IDS2017 was introduced by Sharafaldin et al.\ \cite{sharafaldin2018}. Engelen et al.\ report problems in it: ``TCP appendix'' artefact flows (about a quarter of all flows), attack attempts labelled as attacks although they contain no malicious payload, and feature-extraction flaws that enable shortcut learning; they released a corrected version \cite{engelen2021}. Panigrahi and Borah likewise describe shortcomings of the dataset and propose a combined dataset \cite{panigrahi2018}. We use the \emph{original} labelled flows, so absolute scores here may partly reflect such artefacts.

\textit{Dimensionality reduction for NIDS.} Abdulhammed et al.\ apply an auto-encoder and PCA to CIC-IDS2017, reducing the feature count from 81 to 10 while reporting 99.6\% accuracy with random forest and other classifiers \cite{abdulhammed2019}. Kuang et al.\ propose a hybrid of KPCA, an SVM and a genetic algorithm for intrusion detection \cite{kuang2014}. This section is not a systematic survey; it positions the present study rather than claiming to cover the literature.

\textit{Evaluation practice.} Arp et al.\ describe methodological pitfalls in machine learning for security that can lead to unrealistic performance \cite{arp2022}, and Apruzzese et al.\ propose a framework for cross-evaluating ML-based NIDS across several datasets and discuss both its potential and its risks \cite{apruzzese2022}. Motivated by these works, we (a) fit every data-dependent step on training folds only, (b) report class-aware metrics next to accuracy, (c) include trivial and no-reduction baselines, and (d) state plainly what our sample cannot show.

\section{Dataset and Preprocessing}\label{sec:data}
\textbf{Data.} The study uses the cached 1{,}000-record stratified sample shipped with the project (\texttt{data/sample.npz}), drawn from the Tuesday, Wednesday and Thursday-morning files of CIC-IDS2017 (1{,}308{,}978 flows according to the project documentation; the raw files were not available in the authoring environment, so this count was not recomputed). The sample has 78 numeric features, 11 classes and at least 10 records per class (Table~\ref{tab:classes}); 718 records (71.8\%) are benign. Ten features are constant in the sample. The sample was drawn by the project's \texttt{prepare\_data.py}.

\textbf{A leakage caveat in the cached sample.} The project's original script mean-imputed missing values using all flows \emph{before} sampling. The imputed values in the cached sample therefore carry information from the full dataset. We could not undo this without the raw files; the pipeline below still performs imputation inside each training fold (a no-op on this sample), and the raw-data script \texttt{prepare\_raw.py} removes the issue for future runs. The practical effect is likely small (a mean over 1.3\,M flows) but is not measured.

\begin{table}[t]
\caption{Class distribution of the 1{,}000-record sample.}\label{tab:classes}
\centering\tabsize
\input{tables/T2_class_distribution.tex}
\end{table}

\section{Methodology}\label{sec:method}
\textbf{Pipeline.} Fig.~\ref{fig:pipe} shows the protocol. Within each training fold we fit a mean imputer and a standard scaler, then the dimensionality reducer; test folds are only transformed. Feature spaces are: \emph{original} (78 standardised features); \emph{PCA}; and \emph{KPCA} with linear, RBF (scikit-learn default $\gamma=1/78$) and cosine kernels \cite{scholkopf1998}, each with $n\in\{5,10,15,20,30\}$ components. The original design tied kernel and component count (RBF--5, linear--10, cosine--15); here they are crossed so the two effects can be separated. The reducer is fitted once per fold and shared by all classifiers, which is equivalent to fitting it inside each pipeline because it is deterministic (verified by a unit test against a scikit-learn \texttt{Pipeline}).

\begin{figure*}[t]\centering
\includegraphics[width=\textwidth]{fig1_pipeline.pdf}
\caption{Experimental pipeline. All data-dependent preprocessing and reduction is fitted on the training fold only.}\label{fig:pipe}
\end{figure*}

\textbf{Classifiers.} Logistic regression (L2, lbfgs), $k$-NN with $k=5$ \cite{cover1967}, an SVM with RBF kernel, $C=1$ and $\gamma$=``scale'' \cite{cortes1995}, a decision tree with entropy criterion \cite{quinlan1986}, and a random forest with 100 trees \cite{breiman2001}, using scikit-learn \cite{pedregosa2011}. Hyper-parameters are fixed to the project's original values and \emph{not} tuned. Each classifier except $k$-NN is also run with \texttt{class\_weight=balanced}. A majority-class classifier is the trivial baseline. Resampling (e.g.\ SMOTE) is not used: some classes have only 7--8 training records per fold, too few for meaningful synthetic interpolation.

\textbf{Validation.} Repeated stratified 5-fold cross-validation, 3 repeats (15 folds), seed 42; every method sees identical folds, giving paired observations. Rare classes have 2--3 test records per fold, so per-fold macro-F1 is noisy; we therefore report means and standard deviations over the 15 folds.

\textbf{Nested cross-validation for selection.} Choosing the best of 189 configurations (21 feature spaces $\times$ 9 classifier/weighting variants) on the folds that also produce the reported score is optimistic. We therefore added a nested evaluation on the \emph{same} 15 outer folds: inside each outer training set (800 records), a stratified 3-fold inner cross-validation scores every configuration by macro-F1, the best one is selected, and only that configuration is scored on the untouched outer test fold. We apply it to (A) the global selection, and per classifier to (C) the best \emph{reduced} space and weighting versus (D) the original space with the weighting chosen the same way. Inner validation folds hold only 2--3 records of the smallest classes, so selection is noisy; reducers, scalers and classifiers are refitted on each inner training split. A hold-out test set was not feasible because the smallest classes have 10 records.

\textbf{Metrics.} Accuracy; macro-precision, macro-recall and macro-F1; weighted-F1; balanced accuracy; and the multi-class Matthews correlation coefficient (MCC) as implemented in scikit-learn (Chicco and Jurman argue for MCC in the binary case \cite{chicco2020}). \textbf{Macro-F1 is the primary model-selection criterion}, chosen before the runs because it weights all 11 classes equally. Per-class results and the confusion matrix use pooled out-of-fold predictions of the first repeat.

\textbf{Statistics.} For each classifier we compare feature spaces at a matched $n=10$ with two-sided Wilcoxon signed-rank tests on the 15 paired macro-F1 values, Holm-corrected over the 30 comparisons, with the matched-pairs rank-biserial correlation $r$ as effect size, and Friedman omnibus tests, following the non-parametric procedures discussed by Dem\v{s}ar \cite{demsar2006} (written for comparisons over multiple data sets; here the blocks are cross-validation folds of a single sample). Because cross-validation folds share training data, the 15 differences are not independent and the $p$-values are optimistic \cite{nadeau2003}; we use them for screening, not as proof. A second family compares the original space with each classifier's best reduced space (selected on the same folds, which favours the reduced space); the nested analysis above removes that bias.

\textbf{Cost.} A serial benchmark (no parallel workers) records training and inference time (median of 5 repetitions, one fold split) and peak memory from Python's \texttt{tracemalloc}, measured in a separate pass. Native allocations inside compiled code are not traced, so memory figures are lower bounds that are most reliable for comparing the reducers.

\section{Experimental Setup}
All experiments ran on a 2-vCPU Intel Xeon @ 2.10\,GHz virtual machine with 7\,GB RAM (Linux, Python 3.13.16, scikit-learn 1.9.1, NumPy 2.5.3, SciPy 1.18.1). The main grid contains 21 feature spaces $\times$ 5 classifiers $\times$ 2 weighting schemes (plus the baseline) $\times$ 15 folds = 2{,}850 fits, run with two parallel workers; timings in Table~\ref{tab:rt} come from the separate serial benchmark. A second run repeats the five main spaces at $n=10$ with 10 random stratified splits at test fractions 0.2, 0.4 and 0.6 (training size 800, 600, 400) to examine sensitivity to training-set size. The nested cross-validation (189 candidates $\times$ 15 outer folds $\times$ 3 inner folds, plus outer scoring) took about 6.4 minutes on the same machine. This is \emph{not} a temporal or cross-dataset split: the cached sample has no timestamps or day labels.

\section{Results}
\subsection{Original-feature baseline and the cost of reading accuracy alone}
Predicting the majority class gives accuracy 0.718 but macro-F1 0.076 and MCC 0.000. The random forest on original features reaches accuracy $0.952\pm0.010$ but macro-F1 $0.722\pm0.058$ (weighted-F1 0.947, balanced accuracy 0.717, MCC $0.894\pm0.022$): the gap between accuracy and macro-F1 comes from rare classes. Per class (Table~\ref{tab:perclass}), benign, DoS Hulk, FTP-Patator, SSH-Patator and Heartbleed are detected with F1 $\ge$ 0.96, whereas the three web attacks (brute force, SQL injection, XSS) obtain F1 between 0.21 and 0.50 on only 10--11 records each, and slow-DoS recalls are 0.57--0.71. These are the weak points that accuracy hides (Figs.~\ref{fig:acc} and~\ref{fig:f1}).

\begin{table}[t]
\caption{Per-class results of the selected model (RF, original features), pooled out-of-fold predictions, repeat 1.}\label{tab:perclass}
\centering\tabsize
\input{tables/T8_per_class_best.tex}
\end{table}

\begin{figure*}[t]\centering
\includegraphics[width=\textwidth]{fig3_accuracy.pdf}\\[2pt]
\includegraphics[width=\textwidth]{fig4_macro_f1.pdf}
\caption{Accuracy (top) and macro-F1 (bottom), mean $\pm$ std over 15 folds, 10 components for reduced spaces, standard class weights. Dashed line: majority-class classifier. Axes start at zero.}\label{fig:acc}\label{fig:f1}
\end{figure*}

\subsection{Original space versus PCA versus Kernel PCA}
Table~\ref{tab:n10} and Fig.~\ref{fig:f1} compare spaces at $n=10$. For every classifier, the original space has the highest or (for $k$-NN) statistically indistinguishable macro-F1. Of the 30 Holm-corrected comparisons, 17 were significant; in every significant comparison involving the original space it was the better one, and RBF-KPCA was significantly worse than PCA for LR, SVM, DT and RF. Cosine KPCA differed significantly from PCA only for logistic regression (where PCA was better). Against each classifier's \emph{best} reduced space (Table~\ref{tab:best}), the original features still won significantly for LR, DT and RF (Holm $p<0.01$, large effects); $k$-NN showed no difference; and the SVM was the single clear case where reduction helped: cosine KPCA with 30 components raised macro-F1 from 0.388 to 0.510 ($p_\mathrm{Holm}=0.0003$, better in all 15 folds). That comparison is selected on the same folds and therefore optimistic, and the SVM baseline itself is weak (class weighting alone lifts it to 0.546), so we read it as a sign that default-$\gamma$ RBF-SVM is poorly matched to heavy-tailed standardised features rather than as a general benefit of KPCA.

\begin{table*}[t]
\caption{Accuracy / macro-F1 ($\pm$ std over 15 folds) at 10 components (original = 78 features), standard class weights.}\label{tab:n10}
\centering\tabsize
\input{tables/T7_cv_summary_n10.tex}
\end{table*}

\begin{table}[t]
\caption{Original features versus each classifier's best reduced space (Wilcoxon, Holm-corrected). Positive $\Delta$: original better.}\label{tab:best}
\centering\tabsize
\resizebox{\columnwidth}{!}{\input{tables/T10b_stats_best_reduced.tex}}
\end{table}

\subsection{Selection optimism: nested cross-validation}\label{sec:results-nested}
Selecting the best of 189 configurations on the scoring folds is optimistic, so we repeated the selection inside a nested cross-validation (Section~\ref{sec:method}). The configuration chosen on inner folds reached macro-F1 $0.715\pm0.044$ (accuracy 0.945, MCC 0.880) on the outer folds, against 0.722 for the same-CV maximum: an optimism of about 0.008, small relative to the fold standard deviations. The inner cross-validation chose an original-feature configuration in all 15 outer folds (random forest 6 times, decision tree with standard or balanced weights 5 and 4 times); these three are statistically hard to separate here, and the inner validation folds are very small, so the choice among them should not be over-read. Table~\ref{tab:nested} compares, per classifier, the original space with the best reduced space and weighting chosen in the inner folds, so that neither side benefits from best-of-many selection on its own test data. Original features remained clearly better for DT ($+0.077$) and RF ($+0.071$), only modestly and not significantly at 0.05 for LR ($+0.019$, $p=0.055$, smaller than the same-CV gap of 0.043 in Table~\ref{tab:best}), and indistinguishable for $k$-NN ($+0.010$, $p=0.28$). The RBF-SVM benefitted from reduction (0.595 versus 0.546; the inner folds chose cosine KPCA with 30 components and balanced weights in all 15 outer folds), so the SVM exception survives nested selection but depends on the combination of cosine KPCA and class weighting. The $p$-values are unadjusted paired Wilcoxon tests over 15 outer folds that share training data (Section~\ref{sec:method}); with Holm correction over the five classifiers the DT, RF and SVM conclusions are unchanged and LR stays non-significant.

\begin{table}[t]
\caption{Nested-CV macro-F1 ($\pm$ std over 15 outer folds): original features with the class weighting chosen on inner folds (D) versus the best reduced space and weighting chosen on inner folds (C). $\Delta=$ D$-$C; $p$: unadjusted paired Wilcoxon.}\label{tab:nested}
\centering\tabsize
\resizebox{\columnwidth}{!}{\input{tables/T11_nested.tex}}
\end{table}

\subsection{Kernel and component count}
Linear KPCA is mathematically equivalent to PCA on centred data; the two were identical for $k$-NN and SVM and within 0.001 for logistic regression; for the tree models the macro-F1 differences were larger but small (decision tree $\le0.011$, random forest $\le0.015$ over all component counts and weightings; Fig.~\ref{fig:comp}), which we attribute to small numerical differences between eigen-solvers changing split choices (not tested). Hence the original project's headline configuration (random forest, ``linear KPCA, 10 components'') is a PCA result, not evidence for a nonlinear kernel. Component count mattered mainly for logistic regression (macro-F1 0.28 at 5 components to 0.57 at 30 with PCA, still below 0.61 on the original features) and for the cosine-KPCA SVM; for $k$-NN, DT and RF the curves are flat (Fig.~\ref{fig:comp}). RBF-KPCA with the default $\gamma$ was the weakest kernel for most classifiers; we did not tune $\gamma$, so this is a statement about the default, not about RBF-KPCA in general.

\begin{figure*}[t]\centering
\includegraphics[width=\textwidth]{fig6_components_vs_f1.pdf}
\caption{Macro-F1 versus number of components for PCA and three KPCA kernels; dashed line: original 78 features.}\label{fig:comp}
\end{figure*}

\subsection{Class weighting}
Balanced class weights raised balanced accuracy for LR (0.638 $\to$ 0.778) and SVM (0.358 $\to$ 0.720) on the original features, but macro-F1 improved only for LR (0.612 $\to$ 0.630) and SVM (0.388 $\to$ 0.546) and \emph{fell} for DT and RF (0.713 $\to$ 0.702, 0.722 $\to$ 0.687). The best macro-F1 overall nonetheless remained the unweighted random forest, so the selected model does not depend on this choice.

\subsection{Computational cost}
Table~\ref{tab:rt} summarises the serial benchmark at 800 training / 200 test records. KPCA fitting took about 0.09--0.12\,s against 0.01\,s for PCA, increased decision-tree training time from 0.016\,s to roughly 0.09--0.10\,s and random-forest training from 0.21\,s to 0.36--0.38\,s, and roughly doubled random-forest inference time (0.007\,s to 0.014--0.016\,s per 200 records), since each new sample needs kernel evaluations against all training records. Peak traced memory was 11.7\,MB for KPCA against 1.7\,MB for PCA or no reduction. Over $n=400,600,800,1000$ the traced peak memory of KPCA grew roughly quadratically (3.1, 6.4, 11.0, 16.8\,MB; log--log slope 1.9 between $n=400$ and $1000$), as expected for an $n\times n$ kernel matrix; single timings of the fit were too noisy to support a scaling exponent and we do not claim one. As a purely \emph{theoretical order-of-magnitude estimate} (not a measurement and not a tested requirement), one dense float64 kernel matrix for all 1{,}308{,}978 flows would occupy $n^2\times 8$ bytes $\approx$ 14\,TB, before any eigen-solver workspace, so exact KPCA on the full dataset would require approximation or sub-sampling (not tested here). Absolute times are sub-second on a shared 2-vCPU machine; differences below roughly a factor of two should not be over-interpreted. We make no real-time claim: only batch latencies on this machine were measured.

\begin{table}[t]
\caption{Serial runtime and traced peak memory (800 train / 200 test), median of 5.}\label{tab:rt}
\centering\tabsize
\resizebox{\columnwidth}{!}{\input{tables/T9c_runtime_compact.tex}}
\end{table}

\subsection{Sensitivity to training-set size}
Fig.~\ref{fig:size} shows macro-F1 at 400, 600 and 800 training records over 10 random splits each. For LR, SVM, DT and RF the original space had the highest mean macro-F1 at every size, so the ranking in Table~\ref{tab:n10} is not an artefact of one training size; for $k$-NN the spaces overlap within one standard deviation. This is a stability check within the same sample and distribution, not a test of generalization to unseen traffic.

\begin{figure*}[t]\centering
\includegraphics[width=\textwidth]{fig9_training_size.pdf}
\caption{Macro-F1 versus training-set size (random stratified splits, 10 per size, mean $\pm$ std); 10 components for reduced spaces.}\label{fig:size}
\end{figure*}

\section{Discussion}
\textbf{Answers to the research questions.} RQ1/RQ2: on this sample, reducing 78 standardised features to 5--30 PCA or KPCA components did not improve and usually lowered macro-F1; the one robust gain was for the RBF-SVM with cosine KPCA. RQ3: for $k$-NN, DT and RF the kernel moved macro-F1 slightly more than the component count (both effects small), whereas for LR and SVM component count mattered at least as much; RBF (default $\gamma$) was worst for most classifiers, and linear KPCA equals PCA. RQ4: KPCA multiplied fitting time about tenfold and memory about sevenfold at $n=800$ and scales quadratically, with no accuracy benefit to offset it for four of five classifiers. RQ5: orderings were stable across repeated folds and training sizes within the sample.

\textbf{Why accuracy misleads.} The benign class makes up 71.8\% of records, so the trivial baseline already has accuracy 0.718; high accuracies of 0.93--0.95 for the best models correspond to macro-F1 of 0.66--0.72. The original project reported weighted-F1 and accuracy from single splits, which hide rare-class failure and cannot show whether a difference of a point or two is noise (the standard deviation of macro-F1 across folds is 0.05--0.07).

\textbf{Plausible but untested explanations.} Tree ensembles are insensitive to monotone feature scaling and can exploit individual informative features; projecting onto a few leading components may discard low-variance but discriminative directions. Heavy-tailed flow features also make variance-based reduction and the default RBF width poorly matched. These are hypotheses we did not test (e.g.\ with tuned $\gamma$, robust scaling, or supervised reduction).

\section{Threats to Validity and Limitations}
These limitations are substantial, and the conclusions should be read as conditional on them.

\textbf{1. Very small sample.} All results come from \emph{one} stratified sample of 1{,}000 flow records out of roughly 1.3\,M in the three files used. Nine of the ten attack classes have only 10--17 records, and each cross-validation test fold holds only 2--3 records of the smallest classes. Per-class estimates and macro-F1 are therefore very uncertain (fold standard deviations 0.04--0.07), and the web-attack results in particular rest on 10--11 records per class. A different 1{,}000-record sample could change rankings that differ by a few points.

\textbf{2. Residual leakage in the cached sample.} The project's original script mean-imputed missing values using \emph{all} flows before the sample was drawn (Section~\ref{sec:data}). Scaling, reduction and (formally) imputation are fitted on training folds in our pipeline, but the imputed values in the cached sample already carry full-data information, and we could not remove this without the raw files. We therefore do \emph{not} describe the evaluation as leakage-free. The effect is probably small (it is a mean over 1.3\,M flows, used only to fill missing or infinite values; the share of such values was not quantified) but was not measured; \texttt{prepare\_raw.py} removes it for future runs, and has not been executed here.

\textbf{3. No external validation.} We performed no temporal, cross-day, cross-dataset or held-out-test evaluation \cite{apruzzese2022,arp2022}. Random stratified splits of one sample from one capture may overstate performance, for example because flows from the same attack session can fall into both training and test folds (not measured here). A day-wise split was not possible with the cached sample (no timestamps or day labels), and the provided \texttt{run\_temporal.py} has not been run on real data. Conclusions apply to random splits of this sample only.

\textbf{4. Selection and inference.} The headline configuration (RF, original features) was the best of 189 on the same cross-validation that produced its score, so that score is optimistic. The nested analysis (Section~\ref{sec:results-nested}) estimates the optimism at about 0.008 macro-F1 for the global selection, but it is itself noisy (2--3 records per rare class in the inner folds) and trains its inner models on about 533 records, whereas the outer models train on 800. Repeated-CV significance tests violate independence \cite{nadeau2003}; $p$-values are screening evidence, not proof, and unadjusted values are reported for the nested comparisons.

\textbf{5. No tuning.} Hyper-parameters, including the RBF width $\gamma$, were fixed to the project's original values. Tuning could change the SVM result in particular, and RBF-KPCA conclusions are about the default $\gamma$ only. Class weighting was the only imbalance treatment; resampling methods were not evaluated.

\textbf{6. Dataset quality and scope.} The original CIC-IDS2017 flows contain known labelling and feature-construction artefacts \cite{engelen2021}, which may inflate or distort scores; we did not use the corrected version. Only three of the eight daily files were sampled, so some attack families are absent, and each included family occurs on a single day.

\textbf{7. Cost measurements.} Timings and memory come from one 2-vCPU virtual machine, and memory is a traced Python/NumPy lower bound. The 14\,TB figure is arithmetic ($n^2\times8$ bytes), not a measurement. No approximate KPCA method was tested.

We therefore make no claim of state-of-the-art, real-time, deployable or generalizable performance. The supported finding is narrower: on this sample, under these default settings, reduction by PCA or KPCA did not generally help and sometimes hurt, and accuracy badly overstated detection of rare attack classes.

\section{Conclusion}
On a 1{,}000-record sample of CIC-IDS2017 and under repeated cross-validation with fold-wise fitting of scaling and reduction, a random forest on the original 78 features had the highest observed macro-F1 among the evaluated configurations (macro-F1 $0.722\pm0.058$, accuracy $0.952\pm0.010$; $0.715\pm0.044$ when the selection itself is nested). The 0.23 gap between accuracy and macro-F1 is a central finding: accuracy hides weak detection of web attacks. PCA and Kernel PCA with 5--30 components generally did not improve on the original features for decision tree and random forest (clearly) and logistic regression (modestly), made no meaningful difference for $k$-NN, and helped mainly the RBF-SVM (cosine KPCA, 30 components, with class weighting), at a roughly tenfold fitting-time and sevenfold memory cost with quadratic scaling. Linear KPCA coincided with PCA, so previously reported ``KPCA'' gains with a linear kernel are linear-PCA gains. Accuracy alone concealed weak detection of web attacks. The results are limited to this sample and to random splits.

\textbf{Future work.} Rerun with the full or larger CIC-IDS2017 subsets using the provided \texttt{prepare\_raw.py}; add chronological within-day and cross-dataset evaluation; tune $\gamma$ and use robust scaling; test the corrected CIC-IDS2017 release \cite{engelen2021}; and evaluate approximate KPCA for scalability.

\section*{Reproducibility}
Code, per-fold raw results and instructions are in the project repository (\texttt{RUN\_EXPERIMENTS.md}); the main grid runs in a few minutes on two cores. Data come from CIC-IDS2017 \cite{sharafaldin2018}.

\input{refs_generated.tex}

\end{document}
